% !TEX root = ../../main.tex

\subsection{Sibling Test: Are the Two Children Distinct?}

The sibling test is the second inferential step. For each parent $u$ with
exactly two children $l(u)$ and $r(u)$, it asks whether the two children are
genuinely different from one another rather than merely both differing from the
parent. The edge test alone is not enough to justify keeping a split. A parent
can contain one deviating branch, or a weak left-right asymmetry, without
supporting two stable child clusters. The sibling test supplies that missing
question: once a branch has left its ancestral context, are the two children
actually distinct enough to become separate clusters?
Because the hierarchy is a rooted binary tree, this sibling relation has a
fixed form: every non-root node has exactly one parent and therefore at most one
sibling, namely the other child of that same parent. The sibling test is thus
always a single pairwise comparison, not a comparison against a larger sibling
set.

\paragraph{Sibling null and pooled baseline.}
Formally, it tests the sibling null
\begin{equation}
H_0^{\mathrm{sib}}(u): \theta_{l(u)} = \theta_{r(u)}.
\end{equation}
For binary or indicator coordinates, the pooled null proportion is
\begin{equation}
\bar{\theta}_{u,j}
=
\frac{n_{l(u)}\theta_{l(u),j} + n_{r(u)}\theta_{r(u),j}}
{n_{l(u)} + n_{r(u)}}.
\end{equation}

\paragraph{Standardized sibling contrast.}
Because the two children are disjoint subtrees under the same parent, the
coordinatewise variance takes the usual two-sample Bernoulli form rather than
the nested child-parent form used in the edge test. The standardized sibling
vector is
\begin{equation}
z_{u,j}^{\mathrm{sib}}
=
\frac{\theta_{l(u),j} - \theta_{r(u),j}}
{\sqrt{
\bar{\theta}_{u,j}(1-\bar{\theta}_{u,j})
\left(\frac{1}{n_{l(u)}} + \frac{1}{n_{r(u)}}\right)
}}.
\end{equation}
The sibling test then asks how to combine this multi-feature contrast into one
summary score. The method reuses the spectral information already
computed during the child-parent stage rather than drawing a fresh random
projection from scratch. For a binary parent $u$ with children $l(u)$ and
$r(u)$, let $k_{l(u)}^{\mathrm{edge}}$ and $k_{r(u)}^{\mathrm{edge}}$ denote
the positive edge-stage spectral dimensions of the two children. The sibling
dimension used in code is
\begin{equation}
k_u^{\mathrm{sib}}
=
\operatorname{round}
\!\left(
\sqrt{
k_{l(u)}^{\mathrm{edge}}
k_{r(u)}^{\mathrm{edge}}
}
\right).
\end{equation}
If exactly one child has a positive edge-stage spectral dimension, that value
is reused directly; if neither child has a positive edge-derived dimension, the
parent's own edge-stage spectral dimension is used. The sibling test is run only
for binary sibling pairs with complete Gate 2 spectral context. This standardized
vector answers the first sibling question: how large is the left-right contrast after accounting for
feature-specific sampling noise? The next question is how to combine those
coordinatewise contrasts into one summary score.

\paragraph{Projected sibling statistic.}
The sibling test then combines this multi-feature contrast into a single
summary score while preserving the low-dimensional structure already learned at
the parent during the edge stage.

The projection basis starts with the parent node, not a fresh random draw. The
first available parent PCA directions, where PCA refers to principal-component
directions summarizing the main variation at that parent, are used up to
$k_u^{\mathrm{sib}}$. The requested sibling dimension is capped to the available
parent PCA rank, so the projection basis and eigenvalue vector have the same
dimension.

If $P_u \in \mathbb{R}^{k_u^{\mathrm{sib}} \times p}$ denotes this sibling
projection basis, the raw sibling statistic is
\begin{equation}
T_u^{\mathrm{raw}}
=
\|P_u z_u^{\mathrm{sib}}\|_2^2.
\end{equation}

\paragraph{Reference-law calibration.}
The default p-value calculation is deliberately conservative. After
projection, the method asks whether the observed squared contrast is larger
than would be expected from ordinary within-parent sampling noise, while still
allowing different projected directions to have different variances. In the
method described here, this is done with a Satterthwaite moment match, meaning
an approximation that replaces the exact null law by a scaled chi-square law
with the same first two moments, for the
PCA-derived part of the test and plain $\chi^2(1)$ contributions for any added
orthogonal-complement directions.

\paragraph{Worked sibling example.}
A small numerical example helps separate three ideas that can be conflated:
projection, per-component whitening, and Satterthwaite calibration. Suppose
the standardized sibling difference vector is
\begin{equation}
z_u^{\mathrm{sib}} =
\begin{bmatrix}
2.0\\
-1.0\\
0.5\\
1.0
\end{bmatrix},
\qquad
P_u =
\begin{bmatrix}
1&0&0&0\\
0&1&0&0\\
0&0&1&0
\end{bmatrix}.
\end{equation}
Then the projected contrast is
\begin{equation}
P_u z_u^{\mathrm{sib}} =
\begin{bmatrix}
2.0\\
-1.0\\
0.5
\end{bmatrix}.
\end{equation}
Without any PCA variance correction, the raw quadratic form is
\begin{equation}
T_u^{\mathrm{raw}}
=
\|P_u z_u^{\mathrm{sib}}\|_2^2
=
2.0^2 + (-1.0)^2 + 0.5^2
=
5.25.
\end{equation}
If those three projected coordinates were all unit-variance Gaussian noise,
this statistic would be compared directly to $\chi^2_3$. Now suppose the
first two projected directions are parent PCA directions with eigenvalues
$\lambda_1 = 4$ and $\lambda_2 = 1$, while the third coordinate is an added
unit-variance direction. In matrix form,
\begin{equation}
\Lambda =
\begin{bmatrix}
4&0\\
0&1
\end{bmatrix},
\qquad
\Lambda^{-1} =
\begin{bmatrix}
1/4&0\\
0&1
\end{bmatrix}.
\end{equation}
Per-component whitening would then use
\begin{equation}
T_u^{\mathrm{white}}
=
\begin{bmatrix}
2.0 & -1.0
\end{bmatrix}
\Lambda^{-1}
\begin{bmatrix}
2.0\\
-1.0
\end{bmatrix}
+
0.5^2
=
\frac{2.0^2}{4}
+
\frac{(-1.0)^2}{1}
+
0.5^2
=
2.25.
\end{equation}
The first direction contributes less after whitening because a contrast of
size $2.0$ is less surprising along a direction whose null variance is already
$4$. The implemented sibling default instead keeps the raw quadratic form
$T_u^{\mathrm{raw}} = 5.25$ and changes the reference law. The PCA part is a
weighted sum of $\chi^2(1)$ terms with weights $4$ and $1$, and the added
direction contributes one ordinary $\chi^2(1)$ term. Hence
\begin{equation}
\mathbb{E}[T_u^{\mathrm{raw}}] = 4 + 1 + 1 = 6,
\qquad
\mathrm{Var}(T_u^{\mathrm{raw}}) = 2(4^2 + 1^2 + 1^2) = 36.
\end{equation}
Matching those moments to $c\chi^2(\nu)$ gives
\begin{equation}
c = \frac{36}{2 \cdot 6} = 3,
\qquad
\nu = \frac{2 \cdot 6^2}{36} = 2,
\end{equation}
so the calibration becomes
\begin{equation}
T_u^{\mathrm{raw}} \approx 3\chi^2(2),
\qquad
p_u^{\mathrm{sib}}
=
\Pr\!\left(
\chi^2_2 \ge \frac{5.25}{3}
\right).
\end{equation}
In general, the projected-Wald kernel records this unselected reference as
\begin{equation}
T_u^{\mathrm{raw}}
\approx
a_u \chi^2_{\nu_u},
\qquad
a_u =
\frac{\sum_i \lambda_{u,i}^2}{\sum_i \lambda_{u,i}},
\qquad
\nu_u =
\frac{(\sum_i \lambda_{u,i})^2}{\sum_i \lambda_{u,i}^2}.
\end{equation}
This is the practical difference between whitening and Satterthwaite
calibration: whitening rescales the coordinates before summing, whereas
Satterthwaite keeps the raw quadratic form and adjusts the reference
distribution instead. \Cref{app:projected-wald-example} expands this
example in full matrix form and shows the corresponding plain, whitened, and
Satterthwaite-calibrated calculations side by side.

\paragraph{Post-selection scale distortion.}
Raw sibling $p$-values are still not enough on their own because the sibling
pair being tested is not fixed in advance. The same data are used first to
construct the hierarchy and thereby expose pairs of child subtrees that already
look well separated, and then again to test whether those selected siblings are
significantly different. This data reuse creates \emph{post-selection
scale distortion}. Even under a null setting with no true cluster structure, the tree
construction step tends to present unusually separated sibling pairs because it
was built from the observed dissimilarities. The sibling statistic
is therefore evaluated conditionally on a favorable selection event rather than
under the unconditional null distribution for a pre-specified pair. The
practical consequence is anti-conservative inference: the observed sibling
difference vector $z_u^{\mathrm{sib}}$ is systematically larger than the naive
reference law expects, so raw $p$-values are too small and false split
decisions become too common. Empirical null simulations confirm this effect:
without correction, the Wald sibling test produces Type~I error rates of
approximately 39\% at $\alpha = 0.05$.

\paragraph{Context-weighted empirical-null scale.}
The sibling calibration corrects for this extra optimism by estimating the
runtime empirical-null scale of the projected-Wald statistic from the empirical
null visible inside the fitted hierarchy. For every binary parent $q$, let
$T_q^{\mathrm{raw}}$ be the raw sibling statistic, let $a_q$ and $\nu_q$ be the
Satterthwaite reference scale and effective degrees of freedom returned by the
projected-Wald kernel, let
$\pi_q \in [0,1]$ be the sibling null weight implied by the two child-parent
edge tests, and let $z_q$ denote the calibration context of the sibling
comparison. In the current implementation,
\begin{equation}
\pi_q =
p_{q \to l(q)}^{\mathrm{edge,BH}}\,
p_{q \to r(q)}^{\mathrm{edge,BH}},
\end{equation}
so both child-parent edges must look empirically null-like for the sibling pair
to carry high calibration weight. This is a monotone empirical weight, not a
posterior null probability. Also,
\[
z_q = \log s_q,
\]
where $s_q$ is the resolved sibling projection dimension. The method assumes the
post-selection null family
\begin{equation}\label{eq:sibling-context-scale-model}
T_q^{\mathrm{raw}} \mid H_0
\approx
c(z_q)\,\chi^2_{\nu_q}.
\end{equation}
Here $c(z_q)$ is the total empirical runtime scale. It is not forced to equal
the Satterthwaite scale $a_q$ times a separate selection multiplier, because the
same fitted hierarchy changes both the apparent sibling contrast and the useful
finite-sample scale of the projected quadratic form.
For a focal sibling pair $u$, define context weights
\begin{equation}
w_q(u)
=
\pi_q K(z_q,z_u),
\end{equation}
where the default kernel is Gaussian in the scalar context,
\begin{equation}
K(z_q,z_u)
=
\exp\!\left[
-\frac{1}{2}
\left(
\frac{z_q-z_u}{h}
\right)^2
\right],
\end{equation}
and $h$ is the null-weighted empirical standard deviation of the observed
contexts. The local empirical-null scale estimate is the weighted chi-square
scale maximum-likelihood estimator
\begin{equation}\label{eq:sibling-context-scale-estimator}
\hat{c}(z_u)
=
\max\left\{
1,\,
\frac{\sum_q w_q(u)T_q^{\mathrm{raw}}}
{\sum_q w_q(u)\nu_q}
\right\}.
\end{equation}
The denominator is required to have positive effective calibration mass. If no
positive-degree sibling records carry positive null weight, the calibration
contract is not satisfied and the implementation raises an error.

\paragraph{Runtime correction and sibling false-discovery-rate correction.}
For each focal pair $u$, the adjusted runtime test is
\begin{equation}
T_u^{\mathrm{adj}} = \frac{T_u^{\mathrm{raw}}}{\hat{c}(z_u)},
\qquad
p_u^{\mathrm{sib}}
=
\Pr\!\left(\chi^2_{\nu_u} \ge T_u^{\mathrm{adj}}\right).
\end{equation}
This adjusted test should be read as a practical calibration layer for
post-selection scale distortion rather than as a
separate finite-sample exact theorem for the adjusted sibling statistic.

The sibling multiplicity correction is ordinary
Benjamini--Hochberg \citep{benjamini1995controlling} at
$\alpha_{\mathrm{sib}} = 0.01$, applied only across focal sibling pairs. Let
$S^{\mathrm{sib}}(u) \in \{0,1\}$ denote the resulting indicator that parent
$u$ supports a significant sibling split.
